% SOURCE docs/release/2026-10-29/source/docs/writing-repair/2026-10-29/src/ad84e7f88d7e93fe.tex CHAR 9548
\chapter{Servicing decisions change the next decision's information}
\label{rewrite:servicing}

A servicing policy chooses whether to contact a borrower at two dates. Let
$A_0,A_1$ be contact decisions, $L_0$ baseline borrower state and $L_1$ the
state observed between contacts. For this constructed example, $Y$ is
gross contribution in pounds over the two-period horizon, excluding contact
cost. The policy comparison is sustained contact versus no contact, not
the regression coefficient of one contact conditional on every later event.

The exact law developed in Chapter~\ref{r3:longitudinal} is
\begin{align*}
 P(L_0=1)&=0.5,\\
 P(A_0=1\mid L_0)&=0.2+0.6L_0,\\
 P(L_1=1\mid L_0,A_0)&=0.1+0.5L_0+0.2A_0,\\
 P(A_1=1\mid H_1)&=0.15+0.5L_1+0.1A_0,\\
 Y&=20+3A_0+4A_1+10L_1+6L_0.
\end{align*}
Here $H_1=(L_0,A_0,L_1)$ is the history available for the second action.
Independent innovations implement these kernels. The actions have support
for every possible history, and the example has no hidden confounding.
Consequently, replacing both action kernels with a specified policy and
retaining the state-transition law identifies its mean. This is the
g-formula developed from sequential exchangeability in the retained chapter
and its primary treatment. \citep[Chapter 21]{hernan2020}

\section{Average the state produced by each policy}
For a constant policy $a$, the mean intermediate state is
$\E[L_1(a)]=0.1+0.5(0.5)+0.2a=0.35+0.2a$. Hence
\[
 \psi(a)=20+7a+10(0.35+0.2a)+6(0.5)=26.5+9a.
\]
The gross means are \pounds26.50 and \pounds35.50. At a stipulated
\pounds2 per contact, the net contrast is $9-2(2)=\pounds5$.
The reported exact enumeration independently averages the finite state law
and reweights the observed action paths; both recover the same two means.
There is no finite-sample confidence claim in this calculation.

A regression of $Y$ on all four variables correctly recovers action
coefficients 3 and 4. Their sum misses the extra $10(0.2)=2$ contributed
through the change in $L_1$. Removing $L_1$ from the second-stage regression
does not repair the analysis: the historical second action depends on it.
The calculation therefore distinguishes a fitted conditional equation from
the policy mean obtained by averaging that equation over policy-specific
histories. The distinction changes net value from an erroneous \pounds3
to \pounds5 even in this noise-free conditional outcome model.

\section{What changes when the law must be learned}
In an observed study, the kernels and outcome regressions are estimated.
G-computation substitutes fitted transition and outcome laws and averages
under the target policy. Reweighting instead uses ratios of target to
historical action probabilities. Doubly robust constructions combine these
ideas; their longitudinal scores, nuisance restrictions and remainders are
derived in Chapter~\ref{r3:longitudinal}, rather than assumed from their
name. None removes the requirement that the historical records include
the information responsible for action selection.

For sustained contact, a useful diagnostic is the count and weight
distribution of complete histories supporting that policy. Many borrower
rows do not guarantee enough such histories. A policy learned by maximizing
the same estimated value must be evaluated with the additional selection
uncertainty or an honest evaluation sample. An outcome law conditional on
a hidden state introduces another identification problem; a well-fitting
state-space model does not establish that its intervention meaning is correct.

The product implication is to begin with a fixed, short-horizon policy and
observable histories, retaining a direct g-computation and a weighting
comparison. If support is inadequate or an important action determinant is
unrecorded, changing the collection or trial design is more informative than
adding a flexible learner. The positive \pounds5 example explains the
accounting; it does not establish a real contact policy's benefit.

\section{Estimated policies, missing rewards and decision uncertainty}
Chapters~\ref{gap29:servicing} and \ref{gap29:policy-uncertainty} carry the
policy question into fitted data. They use a separate three-decision economy,
with continuous states and rewards already net of the stipulated action cost.
Those reward units are not the pounds in this two-contact example. Whole
customers, rather than isolated rows, supply the sampling units. A separate
training sample chooses the policy; evaluation and resampling hold that
selected policy fixed.

The fitted comparison is informative precisely because it can fail.
Under its maintained state and observation model, improvement coverage is
94\% with an average interval width of 0.711 reward units. Future-dependent
dropout reduces coverage to 6\%; omitting a state predictor reduces it to
61\%. These observed fractions have Monte Carlo uncertainty, reported in
the study. They demonstrate that resampling an estimated model does not
repair an invalid observation assumption.

The financial interpretation follows the reward definition. Where costs are
already subtracted, subtracting contact cost again would double-count it.
Where a bank's reward law and initial distribution must also be estimated,
the supplied-law components of the generated comparison cannot be treated
as known. Outcome follow-up, action support and measured decision history
therefore determine whether its narrower uncertainty can be useful.


